% !PRINTED 517
\weberpage{517}

%<<CONT
which leave the individual rows $A, B \dots S$ in their
place and only permute the indices within each row. These
permutations form a group by themselves, since by the
composition of two of them the rows $A, B \dots
S$ cannot be interchanged. This group $Q$ is a
\gesperrt{normal divisor} of $P$. For if $\varkappa$ is a permutation of
$Q$, and $\pi$ one from $P$, then the composed
permutation $\pi^{-1} \varkappa \pi$ belongs again to $Q$, because the
permutation among the $A, B \dots S$ that is produced by $\pi^{-1}$
is not altered by $\varkappa$ and is undone again by
$\pi$. The group $Q$ is
intransitive, and by adjunction of a function $\psi$ belonging
to $Q$ not only does the \gesperrt{Galois} resolvent
become reducible, but the function $f(x)$ itself splits into $s$
factors of degree $r$,

\weq{2}{f(x) = f_\alpha(x,\psi)\, f_\beta(x,\psi) \dots f_\sigma(x,\psi).}

The roots of $f_\alpha(x,\psi) = 0$ are the quantities $\alpha$,
and the \gesperrt{Galois} group of this equation consists of all the
permutations that are produced by $Q$ among the $\alpha$.

The quantity $\psi$ is, if $j$ is the index of the divisor $Q$ of
$P$, a root of a partial resolvent of degree $j$,
$\chi(u) = 0$, which is a normal equation and hence its
own \gesperrt{Galois} resolvent (\S.\,155). If we denote
its roots by

\weq{3}{\psi,\ \psi_1,\ \psi_2 \dots \psi_{j-1},}

then the \gesperrt{Galois} group of this resolvent consists of the
substitutions

\weq{4}{(\psi,\psi),\ (\psi,\psi_1),\ (\psi,\psi_2) \dots (\psi,\psi_{j-1}),}

which are produced by the permutations of the subsidiary groups

\weq{5}{Q,\ Q\pi_1,\ Q\pi_2 \dots Q\pi_{j-1}}

Each of these subsidiary groups produces a certain
permutation among the rows $A, B \dots S$, and
indeed these permutations are all different from one another;
for if, say, the same permutation among the rows were produced
by $\pi_1$ and $\pi_2$, then $\pi_2 \pi_1^{-1}$ belongs
to $Q$, and $Q\pi_1$, $Q\pi_2$ are not different
subsidiary groups. The permutation group of the partial resolvent
$\chi(u) = 0$ is therefore isomorphic with the group of the
permutations that are produced by $P$ among the rows $A, B \dots S$.

If we understand by $Q_\alpha$ the totality of the permutations
of $P$ by which the row $A$ is not displaced, then
$Q_\alpha$ is
%CONT>>
